\chapter{Java Objects as Facts}
\label{ch:java-objects}

Most \Morendo\ applications reason over Java objects rather than over
\fn{deftemplate} facts. A Java class is declared to the engine, which
turns its bean properties into the slots of a template; instances of the
class are then asserted and matched like any fact. This chapter covers the
CLIPS side and the rules involved. Chapter~\ref{ch:embedding} gives the
Java calls.

\section{Declaring a class}

From Java the call is \fn{engine.declareObject(Account.class)}; from CLIPS
the equivalent is

\begin{lstlisting}
(defclass <classname> [<template-name>] [<parent-template>])
\end{lstlisting}

The engine introspects the class with \fn{java.beans.Introspector}: every
readable property except \fn{class} becomes a slot with the property's
name, typed from the property's Java type, and an array-valued property
becomes a multislot. The template is named after the \emph{fully
qualified class name} unless a template name is given, so a rule over
\fn{woolfel.examples.model.Account} reads

\begin{lstlisting}
(defrule Over_30
    (woolfel.examples.model.Account
        (first ?first)
        (last ?last)
        (age ?age&:(> ?age 30)))
=>
    (printout t ?first " " ?last " is over 30" crlf))
\end{lstlisting}

The example beans live in the \fn{examples} module, in the package
\fn{woolfel.examples.model}; the rule files \file{sample1.clp} and
\file{join\_sample1.clp} under \file{samples/} use them. Declaring a class prints its
template name, which is why transcripts that load these samples begin
with a few class names.

A parent template makes the new template share the parent's slot layout,
so a rule written against the parent also sees instances of the child.
Declaring the same class twice is harmless.

\section{Asserting objects}

Objects are asserted from Java, or from CLIPS with

\begin{lstlisting}
(definstance <object> <template-name>)
\end{lstlisting}

where the object is a bound variable, usually one created with
\fn{(new classname args...)} or returned by a Java function. When an
object is asserted the engine creates a \emph{shadow fact}: a copy of the
property values, given a fact id like any other fact. \fn{(facts)} shows
shadow facts in the template's slot order, and rules match them without
knowing where they came from. Asserting the same object instance twice
does nothing.

\section{Keeping the shadow fact current}

A shadow fact is a snapshot. When the Java object changes, the engine has
to be told, in one of two ways:

\begin{itemize}
\item \textbf{Property change events.} If the class has public
  \fn{addPropertyChangeListener} and \fn{removePropertyChangeListener}
  methods taking a \fn{PropertyChangeListener}, the engine registers
  itself as a listener when the object is asserted. Each
  \fn{PropertyChangeEvent} the object fires makes the engine retract and
  re-assert the shadow fact, so rules see the change at once. The class
  does not have to extend \fn{PropertyChangeSupport}; it only has to fire
  the events, as the \fn{Account} example does in every setter.
\item \textbf{Explicit modification.} Otherwise, after changing the
  object, call \fn{engine.modifyObject(obj)} from Java. Nothing happens
  until then, and a rule that already fired for the old values does not
  fire again by itself.
\end{itemize}

\fn{engine.retractObject(obj)} removes the shadow fact. Objects asserted
as \emph{static} (a flag on the Java \fn{assertObject} call) get no
listener and ignore \fn{modifyObject}; they are meant for reference data
that never changes while the engine runs.

\section{Working with objects in rules}

Inside a rule the slots of a shadow fact are read with the usual pattern
syntax. To call into the object itself, bind the fact and use the Java
functions:

\begin{center}
\begin{tabularx}{\textwidth}{lL}
\toprule
\fn{(get-member ?f property)} & read a property through its getter \\
\fn{(set-member ?f property value)} & write a property through its setter \\
\fn{(call ?f method arg...)}, \fn{(member ...)} & invoke a method by reflection \\
\fn{(new classname arg...)} & construct an object \\
\fn{(instanceof ?f classname)} & class test \\
\fn{(load-package prefix)} & let later calls use short class names \\
\bottomrule
\end{tabularx}
\end{center}

Setting a property from a rule with \fn{set-member} on a bean that fires
change events causes the shadow fact to be modified while the rule is
firing, which is the intended way to update an object from the RHS.
The fact functions of Section~\ref{sec:fact-access} accept a bound
object as well: \fn{(fact-slot-value ?f age)} reads the shadow fact's
slot, \fn{(fact-index ?f)} its id and \fn{(fact-existp ?f)} tells
whether the object is still asserted.

\section{Macros instead of reflection}

Reading properties by reflection costs time on every match. The
\fn{Macro Functions} group generates and loads small compiled accessor
classes for a declared class:

\begin{lstlisting}
(generate-macro woolfel.examples.model.Account)
(use-macro woolfel.examples.model.Account)
\end{lstlisting}

After \fn{use-macro} the engine reads that class's properties through the
generated code. The service module (Chapter~\ref{ch:embedding}) does this
automatically for every class listed in its configuration. The generated
sources and classes go to a directory per module under the one named by
the JVM property \fn{morendo.workdir}, or by default under a
\fn{morendo} directory in the JVM's temporary directory; nothing is
written there until a macro is generated.

\section{Objects that are not beans}

A class with no getters is declared, but its template has no slots to
match. Wrap such objects in a bean, or assert ordinary facts built from
them. The engine also accepts a non-shadow mode in which the fact reads
the live object on every access instead of copying it; it is selected by
the Java \fn{assertObject} call and described in
Chapter~\ref{ch:embedding}.
